% TOP_PROBLEM: {"id": 1346, "title": "Yang-Mills Existence and Mass Gap", "category_id": 6, "category": {"id": 6, "name": "geometry", "display_name": "Geometry", "description": "Euclidean and non-Euclidean geometry, geometric structures.", "slug": "geometry", "order_index": 6, "created_at": "2026-07-31T15:26:25.671Z"}, "status": "open"}
% FIRST_POSTED: null
% ENTRY_KIND: statement_only
% Imported from: batch06.tex; UnsolvedMath ID 1346
\documentclass[11pt]{article}
\usepackage[margin=25mm]{geometry}
\usepackage{fontspec}
\setmainfont{Latin Modern Roman}
\usepackage{amsmath,amssymb,mathtools}
\usepackage{unicode-math}
\setmathfont{Latin Modern Math}
\usepackage{xurl,hyperref}
\hypersetup{colorlinks=true,urlcolor=blue}
\setcounter{secnumdepth}{0}
\setlength{\parindent}{0pt}
\setlength{\parskip}{5pt}
\setlength{\emergencystretch}{3em}
\newcommand{\sref}[2]{\hyperlink{source:#1:#2}{[S#2]}}
\newcommand{\cataloguescope}[1]{\paragraph{Recorded target.}#1}
\begin{document}
% UnsolvedMath ID 1346; source catalogue code GEOM-009
\setcounter{section}{1345}
\section{Yang-Mills Existence and Mass Gap}\label{1346}
{\small\sffamily UnsolvedMath ID 1346 · Geometry}
\subsection{Definitions and mathematical statement}

\paragraph{Shared notation.}$\mathbb N=\{1,2,\ldots\}$, and $\mathbb Z,\mathbb Q,\mathbb R,\mathbb C$ denote the integers, rationals, reals and complex numbers. Any other conventions are specified locally.


Let $G$ be a compact simple Lie gauge group, with Lie algebra $\mathfrak g$ and an invariant positive inner product. A connection $A$ has curvature $F_A=dA+A\wedge A$ and covariant derivative $D_A$. The classical pure Yang--Mills action on Euclidean four-space is proportional to $g^{-2}\int_{\mathbb R^4}|F_A|^2\,dx$, where $g>0$ is the coupling; no matter fields are included.
A quantum theory must have gauge-invariant local operator-valued distributions corresponding, with renormalization, to the local polynomials in $F_A$ and its covariant derivatives. Its correlation functions must have the short-distance behavior of asymptotic freedom and perturbative renormalization, including a local stress tensor and the prescribed operator-product singularities.
In the Minkowski formulation, the theory has a Hilbert space $\mathcal H$, a unitary Poincar\'e representation with energy--momentum spectrum in the closed forward cone, a vacuum $\Omega$ invariant under that representation, and covariant local fields commuting at spacelike separation. Existence is required to satisfy an axiom system at least as strong as the Wightman axioms, or equivalently the Osterwalder--Schrader reconstruction framework, rather than a formal path integral. Let $H$ be the self-adjoint generator of time translations, so $H\Omega=0$ and $\operatorname{Spec}H\subseteq[0,\infty)$. Define
\[
 m=\sup\{\Delta>0:\operatorname{Spec}H\cap(0,\Delta)=\varnothing\}.
\]
\paragraph{Statement recorded in the supplied source.}
For every compact simple gauge group $G$, construct a nontrivial pure quantum Yang--Mills theory on four-dimensional spacetime with the correspondence, short-distance behavior and axiomatic properties specified above, such that its mass satisfies
\[
 0<m<\infty.
\]
Thus the positive-energy spectrum must be separated from zero by a strictly positive gap. Existence on a compact four-torus alone is not the requested existence on $\mathbb R^4$. Confinement and the existence of an isolated one-particle state are separate extensions, not additional requirements of this problem. \sref{1346}{2}
\subsection{Short English statement}
For every compact simple gauge group, construct a mathematically well-defined interacting Yang--Mills quantum field theory in four dimensions and prove that its first positive spectral energy is bounded away from zero.
\subsection{Sources}
\begin{description}
\item[\hypertarget{source:1346:1}{S1}] ULAM.AI and the UnsolvedMath Contributors, \emph{UnsolvedMath: A Curated Collection of Open Mathematics Problems}, record 1346 (GEOM-009). CC BY 4.0. \url{https://huggingface.co/datasets/ulamai/UnsolvedMath}.
\item[\hypertarget{source:1346:2}{S2}] Arthur Jaffe and Edward Witten, \emph{Quantum Yang--Mills Theory}, official Clay Mathematics Institute Millennium Prize problem description, Sections 3--4, pp. 4--6. \url{https://www.claymath.org/wp-content/uploads/2022/06/yangmills.pdf}.
\end{description}
\label{end:1346}
\end{document}
